\scanpage{36}
\[
(18)\qquad \varphi:\ p^*(E_{Y,X})\longrightarrow i^*(E_{X^*}).
\]

In terms of the diagram of groups
$(\unsure{10|1^\circ})$ When the \inserted{three} constituent groupoids of
(8) are connected, i.e. $Y$, $X^*$ and $\mathcal V^*_{Y,X}$
connected), one can also interpret
these systems, \deleted{\unsure{as systems|as a system}}\inserted{by regarding $E_{X^*}$ as}
a set on which $\Pi_1(X^*)$ acts, $E_{Y,X}$
as a set on which $\Pi_1(Y)$ acts,
and $p^*(E_{Y,X})$ and $i^*(E_{X^*})$ in (18) as
being these $=$s sets, \deleted{\unsure{equipped}}\inserted{regarded} as
$\Pi_1(\mathcal V^*_{Y,X})$-sets \deleted{which} by restriction
of the operators.

\margin{NB In the case where one has $=$ a
diagram (14), this means that $\Pi_1(Y)$
appears as a quotient of $D=\Pi_1(\mathcal V^*_{Y,X})$, and
thus $E_{Y,X}$ can be interpreted straightaway as
a $\deleted{\unsure{\Pi_1(\mathcal V^*_{Y,X})}}D$-set, on which the \unsure{subgroup}
$I$ acts trivially, one
interprets $\varphi$ of (18) as a
homomorphism of $\Pi_1(Y)$-sets.
\[
(18\ \mathrm{bis})\qquad E_{Y,X}\longrightarrow
\underbrace{E^I_{X^*}}_{I\text{-invariants}}
\]}

A \inserted{typical,} particularly important case
for \unsure{the|the} \deleted{applications to} Algebraic Geometry,
is that where $X$ is a topological
surface and where $Y$ is reduced to a point $a$,
thus $\mathcal V_{Y,a}=\mathbb D_a$, \uline{disk} \inserted{(infinitesimal)} around $a$,
and $\mathcal V^*_{Y,a}=\mathbb D^*_0$, punctured (infinitesimal) disk around
$a$, so that one finds the diagram
% Pale bottom-left foliation, in another hand; not the author's own numbering.
\margin{35}
